Eight studies provide the scientific foundation for the system's design.
Each addresses a specific aspect of water monitoring, water-heater control,
optical data acquisition, or user interaction.
Together, these studies inform the architectural decisions adopted in this work.

\Needspace{6\baselineskip}
\subsection*{A1. Smart monitoring of water use}
Palermo et al.~\cite{palermo2022} review smart technologies for water resource management,
particularly at building scale. Their overview connects sensing and communication with
information about occupants' water use. This supports the integration of a monitoring
interface that presents observations from the prototype. Providing such information,
however, does not establish water or energy savings for the appliance considered here.

\Needspace{6\baselineskip}
\subsection*{A2. Control of domestic water heating}
Ritchie et al.~\cite{ritchie2022} investigate centrally coordinated control schedules for
multiple electric water heaters. Their simulations compare hot-water delivery strategies
and constrained peak loads. The study motivates distinguishing control actions from
measured outcomes. Unlike their coordinated scheduling approach, the system developed
here interacts with one heater through its existing physical interface, preserving the
original electrical control mechanisms.

\Needspace{6\baselineskip}
\subsection*{A3. Networked water-consumption monitoring}
Slan\'{y} et al.~\cite{slany2020} develop a modular smart-metering platform using
LoRaWAN and evaluate it in laboratory and residential conditions, including leakage detection.
Their work demonstrates the relevance of transmitting water-use measurements for monitoring.
The present system instead uses a Hall-effect turbine flow sensor and Wi-Fi communication.


\Needspace{6\baselineskip}
\subsection*{A4. Flow measurements outside the pipe}
Riboldi et al.~\cite{riboldi2023} investigate contactless monitoring with ultrasonic
transducers and external electrodes clamped onto plastic pipes. Their hydraulic-loop
experiments show an alternative to inserting sensors into the fluid path. This distinguishes
non-invasive access to the heater electronics from non-invasive flow measurement:
our turbine sensor still requires water to pass through its body. Their results do not
validate our sensor or its installation.

\Needspace{6\baselineskip}
\subsection*{A5. Demand and flow data}
Lee and Yim~\cite{lee2021} analyze domestic hot-water energy and flow demand using smart-meter data.
This supports measuring use rather than treating powered time as water consumption.


\Needspace{6\baselineskip}
\subsection*{A6. Optical reading of a display}
Tongbai~\cite{tongbai2024} uses an ESP32-CAM to capture a digital instrument display
without a communication port and recognizes its seven-segment digits on a web server.
We apply the same capture-and-recognize approach to the Rheem visor, subject to testing
its geometry and lighting. Recognition accuracy reported for the calibration instrument
cannot be transferred to the heater display without corresponding tests.

\Needspace{6\baselineskip}
\subsection*{A7. Advanced scheduling as a boundary}
Shen et al.~\cite{shen2021} study data-driven electric water-heater scheduling using
robust model predictive control. Their evaluation combines a real-world dataset with
simulations of a central water heater. Their work explains why demand information may
be useful later.

\Needspace{6\baselineskip}
\subsection*{A8. User-centered smart-home services}
Chang and Nam~\cite{chang2021} examine smart-home service preferences and adoption
through a survey of potential users in South Korea. Their findings associate service
preferences with adoption intentions and identify interest in safety and convenience.
This motivates presenting clear status information and accessible controls. It does not
constitute a usability evaluation of our interface or demonstrate energy savings.

\Needspace{6\baselineskip}
\subsection{Proposed Contribution}
These sources separately address hot-water control, flow data, optical readout, and user interaction.
Our contribution is to combine optical observations of the specific EW1-50RP with removable
button actuation while retaining the heater's original controls, and to demonstrate flow
acquisition separately. The integration addresses monitoring and external interaction;
its evaluation must distinguish observed device states, recognition results, and control
requests rather than infer energy savings from connectivity alone.
